\section{Implementation}
\label{sec:implementation}

MOMO is implemented as a modular control-plane extension over an existing exploratory-QA harness, not as a fork of it. The baseline \texttt{DataLakeAgent} is reused unmodified; MOMO is delivered through a subclass that overrides five generic extension hooks. The same tasks run through the baseline and through each MOMO configuration with no code-path divergence in the baseline itself.

\subsection{Extension Hooks}

\texttt{SanaDataLakeAgent} extends \texttt{DataLakeAgent} through five hooks:

\begin{enumerate}
\item \texttt{\_pre\_build\_setup()} --- resets delegation cost tracking when the \texttt{delegation} flag is on, so per-task subagent cost is isolated.
\item \texttt{\_extra\_prompt\_text()} --- composes prompt blocks for each active flag (\texttt{cot\_block()}, \texttt{delegation\_block()}, \texttt{sprint\_block()}). The blocks document the contract or protocol the model must follow.
\item \texttt{\_extra\_plugins()} --- wires the runtime plugins: \texttt{CoTSteerHandler}, \texttt{SprintSteerHandler}, and a \texttt{StateOfTaskDashboardPlugin} that maintains the persistent state-of-task readout used by \texttt{sprint}.
\item \texttt{\_conversation\_manager()} --- returns a \texttt{SummarizingConversationManager} tuned to preserve plans and render them as completion checklists, helping the planner track \texttt{missing\_outputs} across multiple inspect-contract returns.
\item \texttt{\_decorate\_tools()} --- swaps \texttt{peek\_file} and \texttt{peek\_multiple} for the profile-enriched wrappers when \texttt{results} is on; injects \texttt{search\_subagent} and \texttt{inspect\_subagent} (with bound guards) when \texttt{delegation} is on; injects the \texttt{cot} and \texttt{sprint} tools when those flags are on.
\end{enumerate}

The baseline \texttt{DataLakeAgent} is never modified. With all flags off, \texttt{SanaDataLakeAgent} reduces to the baseline plus an optional summarizing conversation manager. \texttt{SanaBatchRunner} is a one-line subclass of the baseline batch runner: \texttt{\_AGENT\_CLASS = SanaDataLakeAgent}.

\subsection{Flags and Configuration}

Feature toggles live in a \texttt{SanaFlags} dataclass with four opt-in booleans (\texttt{sprint}, \texttt{cot}, \texttt{results}, \texttt{delegation}) plus hyperparameters: the macro-reflection cadence $k$, the \texttt{sprint} mode (\texttt{cadence} or \texttt{commitment}), and per-subagent call budgets. Validation enforces the \texttt{sprint}/\texttt{delegation} mutual-exclusion rule and plan-mode dependencies. \texttt{SanaFlags} is composed into a \texttt{SanaRunConfig} that extends the baseline \texttt{RunConfig}; the rest of the evaluation harness (task loading, tool budgeting, tracing, logging, result writing) is untouched.

\subsection{Delegation Runtime}

The delegation runtime lives in \texttt{sana\_evaluation/tools/delegation\_*.py}. \texttt{DelegationRuntime} owns the contract lifecycle: validate the incoming \texttt{SearchContract} or \texttt{InspectContract} (including budget bounds); spawn a fresh agent with the appropriate tool surface (retrieval tools for search; full data-tool set for inspect); inject the relevant guards (\texttt{\_FileReferenceGuard} for both; \texttt{\_InspectSourceGuard} additionally for inspect); attach a \texttt{\_SubagentBudgetSteer} for budget enforcement and a \texttt{\_SubagentToolLedger} for per-call recording; run the agent until it emits a single \texttt{contract\_return}, until the budget is exhausted, or until a fatal failure; validate the return payload against the schema and surface it to the planner.

The lifecycle makes every contract observable. Each contract has recorded start and end events, budget usage, tools-used list, token usage, and cost. This is the evaluation surface used in Section~\ref{sec:experiments}.
